\scanpage{83}
\mspage{9}
par un tore --- pratiquement, par une ``jacobienne généralisée'' convenable de $X$.

Mais il faut se rappeler que pour une telle variété (p.ex. une courbe elliptique, ou la variété $G_m(\mathbb C)=\mathbb C^*$) l'opération de $\Gamma$ sur le groupoïde $\Pi_1(X_K)$ n'a plus rien de fidèle, donc pour ``repérer'' un $u\in\Gamma$, il ne suffit pas de connaître des invariants du type $g_l(u)$ \uline{relatifs à} \uline{ce même} $X$. Il faudra donc de façon essentielle introduire des ``paramètres de repérage'' pour $u$, en utilisant d'autres variétés algébriques déf sur $K$ --- la plus naturelle étant
\[
\tag{35}\unsure{U|\mathbf U}_{0,3}=\mathbb P^1\setminus\{0,1,\infty\}
\]
\figure{\begin{tikzpicture}\draw(0,0)--(2.8,.12);\end{tikzpicture}}
Faisant maintenant varier $X$, en considérant un morphisme de $K$-schémas
\[
\tag{36}f:X\mathbin{-}\unsure{X|Y}
\]
on déduit des morphismes de groupoïdes
\margin{82}
